\documentclass{article}
\usepackage[utf8]{inputenc}
\usepackage[T1]{fontenc}
\usepackage{enumitem}
\usepackage{hyperref}

\begin{document}

\section*{Onde as áreas se convergem}
\begin{itemize}
    \item \textbf{IA para Gestão de Constelações de Satélites:} Otimizar rotas, consumo de energia e transmissão de dados em redes de satélites (como Starlink ou futuras redes de órbita baixa) é um problema massivo de computação que a IA está resolvendo agora.
    \item \textbf{Comunicações Quânticas via Satélite:} O uso de satélites para distribuição de chaves quânticas (QKD) é a fronteira da segurança de dados global. O MIT tem pesquisas ativas em sistemas fotônicos e comunicações ópticas espaciais.
    \item \textbf{Computação Quântica em Ambientes de Radiação:} Simular como qubits se comportam no vácuo e sob radiação espacial é um nicho técnico de alto valor.
\end{itemize}

\section*{Caminhos Recomendados dentro do MIT}
Para alguém com o seu histórico, os programas que mais fazem sentido são:
\begin{itemize}
    \item \textbf{M.Sc./Ph.D. no Departamento de AeroAstro:} Se você quer focar no hardware/sistema espacial, mas usando IA como ferramenta.
    \item \textbf{M.Sc. em Computer Science (EECS):} Se o seu foco for o desenvolvimento dos algoritmos quânticos ou modelos de IA, aplicando-os posteriormente a problemas do setor espacial.
    \item \textbf{System Design \& Management (SDM):} Este é um mestrado conjunto entre a Escola de Engenharia e a Sloan School of Management. É ideal para quem já tem experiência profissional e quer liderar projetos complexos que unem tecnologia avançada e estratégia de negócio.
\end{itemize}

\section*{Estrutura do Candidato de Elite para o MIT}
Para se destacar em áreas tão técnicas e competitivas, seu planejamento deve seguir este ciclo de preparação:
\begin{itemize}
    \item \textbf{Experiência Prática (Projetos):} Como você já trabalha com software, tente migrar partes da sua atuação profissional para lógica de processamento de sinais ou automação. Se possível, contribua para projetos de código aberto relacionados a SpaceTech ou bibliotecas de Quantum Computing (como a Qiskit da IBM).
    \item \textbf{Publicações e Conferências:} No mundo acadêmico de elite, ser autor ou coautor de artigos em conferências relevantes (mesmo que nacionais ou de nicho em IA) conta muito mais do que apenas ter notas altas.
    \item \textbf{Networking:} Comece a acompanhar os papers publicados pelos professores do MIT CSAIL (Computer Science and Artificial Intelligence Laboratory). Eles são os mesmos que decidem quem entra no mestrado/doutorado.
\end{itemize}

\section*{Dica de Ouro: O "Pulo do Gato"}
Dado o seu interesse em Satélites, pesquise sobre o MIT Space Systems Laboratory (SSL). Eles desenvolvem pequenos satélites (CubeSats) e sistemas de controle. Analise os projetos atuais deles para ver qual tecnologia eles estão usando para automação — isso pode ser o tema central da sua carta de intenção (Statement of Objectives).

\end{document}
